\section{Problem setup}\label{sec:problem-setup}

% TODO: fix the conceptual setup the rest of the paper operates in.  Definitions
% here should be conceptual; concrete constructions belong in \Cref{sec:methods}.

\subsection{Baseline}\label{subsec:baseline}

% TODO: describe the baseline state that this paper takes as fixed (often the
% close of the predecessor paper).  State explicitly which parts are retained
% unchanged throughout this paper.

\subsection{Conceptual objects}\label{subsec:conceptual-objects}

% TODO: introduce the new objects this paper will work with at the conceptual
% level only.  Concrete construction is deferred to \Cref{sec:methods}.

\subsection{Compatibility relations}\label{subsec:compatibility}

% TODO: state how the new objects are required to relate to the inherited
% baseline (e.g., descent compatibility, layer relation, mapping discipline).

\subsection{Boundary conditions}\label{subsec:boundary-conditions}

% TODO: state any irreducibility, non-factorization, or boundary conditions
% required for the paper's later positive claims to be non-trivial.

\subsection{Baseline status ledger}\label{subsec:baseline-status}

% TODO: prose ledger of baseline status; full extended table belongs in the
% extended-material appendix \Cref{app:extended-material}.

\subsection{Conceptual figure}\label{subsec:conceptual-figure}

% TODO: drop in a conceptual figure that summarizes the setup.
% % TEMPLATE: Figure placeholder 1.
% Replace the framed box with the final figure or TikZ/PGF content.
\begin{figure}[tbp]
  \centering
  \fbox{\rule{0pt}{1.5in}\rule{0.8\linewidth}{0pt}}
  \caption{Template figure caption.}
  \label{fig:template-01}
\end{figure}

